% qtes-dftb User Manual -- LaTeX version of docs/MANUAL.md
% Build:  cd docs/manual && make     (or: pdflatex qtes-dftb-manual.tex, twice)
\documentclass[11pt,a4paper]{report}

\usepackage[T1]{fontenc}
\usepackage[utf8]{inputenc}
\usepackage{lmodern}
\usepackage{microtype}
\usepackage[margin=25mm]{geometry}
\usepackage{amsmath}
\usepackage{array}
\usepackage{booktabs}
\usepackage{tabularx}      % short tables: never split
\usepackage{xltabular}     % long tables: may break across pages
\usepackage{needspace}
\usepackage{enumitem}
\usepackage{xcolor}
\usepackage{listings}
\usepackage{caption}
\usepackage[hidelinks]{hyperref}

\setlist{nosep,leftmargin=*}
\setlength{\parindent}{0pt}
\setlength{\parskip}{0.6em}
\renewcommand{\arraystretch}{1.25}
\emergencystretch=3em   % long \texttt names: stretch spaces instead of overflowing

% inline code: file names, variables, commands
\newcommand{\code}[1]{\texttt{\detokenize{#1}}}
% placeholder for content that needs the code authors
\newcommand{\todo}[1][to be written]{\textbf{\textcolor{red!70!black}{[#1]}}}
\newcolumntype{L}[1]{>{\raggedright\arraybackslash}p{#1}}
\newcolumntype{Y}{>{\raggedright\arraybackslash}X}

\definecolor{codebg}{gray}{0.95}
\lstset{
  basicstyle=\ttfamily\small,
  backgroundcolor=\color{codebg},
  frame=none,
  columns=fullflexible,
  keepspaces=true,
  breaklines=true,
  xleftmargin=0.5em,
  xrightmargin=0.5em,
  aboveskip=0.8em,
  belowskip=0.8em,
  upquote=true,
}
\lstdefinestyle{bash}{language=bash,showstringspaces=false}

\title{\textbf{qtes-dftb}\\[0.5em]
  \Large User Manual\\[1em]
  \large Quantum-trajectory dynamics with DFTB forces}
\author{}
\date{Draft, step 1 --- \today}

\begin{document}
\maketitle

\begin{abstract}
\noindent This is a draft. Chapters marked \todo{} need input from the code
authors, mainly the physics background. All other statements were checked
against the source code and the example runs.
\end{abstract}

\tableofcontents

%======================================================================
\chapter{Introduction}

\code{qtes-dftb} simulates the dynamics of a molecular system in which a
few light nuclei (typically one H atom) are treated \textbf{quantum
mechanically} and all other atoms \textbf{classically}. The quantum nuclei
are represented by an ensemble of \textbf{quantum trajectories} (QTs) that
sample their wavepacket. Their quantum effects enter through an approximate
quantum potential, computed with the \textbf{linearized quantum force} (LQF)
method. The forces on all atoms come from \textbf{DFTB} electronic-structure
calculations done on the fly for every trajectory. The classical atoms move
either in the trajectory-averaged (Ehrenfest) force or independently in
each trajectory.

The program is a parallel Fortran code (\code{bohmd}, version 17/08/2016).
It distributes the trajectories over MPI ranks.

\textbf{Typical application:} an H atom scattering from, or chemisorbed on,
a hydrocarbon or a graphene flake. All included examples are of this type.

\section*{Please cite}

\begin{itemize}
\item The method: S.~Garashchuk, J.~Jakowski, L.~Wang, B.~G.~Sumpter
  \cite{garashchuk2013}.
\item The H-on-graphene application: L.~Wang, J.~Jakowski, S.~Garashchuk
  \cite{wang2014}.
\end{itemize}

Related: nonlinear quantum corrections \cite{wang2016}, which are not
implemented in this code version.

%======================================================================
\chapter{Theory in brief}

\todo{} Planned content, to be written with the code authors:
\begin{itemize}
\item the quantum-trajectory (Bohmian) formulation: QT positions, momenta,
  action $s$ and weights $w$;
\item the linearized quantum force (LQF) and the fit used in \code{qpot};
\item the Ehrenfest treatment of the classical atoms;
\item the initial wavepacket and its sampling;
\item the definitions of the energies and of the correlation function
  written to \code{ener} and \code{cor}.
\end{itemize}
These should follow the notation of ref.~\cite{garashchuk2013}.

Facts already established from the code, to be built on:
\begin{itemize}
\item \textbf{Wavepacket:} the initial wavepacket of each quantum coordinate
  is a Gaussian with width parameter $a_0$, centre $q_0$ and momentum $p_0$.
  The sampling spread is
  \begin{equation}
    \sigma = \frac{1}{\sqrt{4 a_0}} .
  \end{equation}
\item \textbf{Quantum-force fit:} the basis is linear, $(x, y, z, 1)$ per
  quantum atom. The fit solves the weighted moment matrix of the QT
  positions.
\item \textbf{Propagation:} velocity Verlet. Each large step $\Delta t$ is
  split into \code{kout} substeps.
\end{itemize}

%======================================================================
\chapter{Installation}

\section{Requirements}

\begin{tabularx}{\linewidth}{L{4.2cm} L{4.2cm} Y}
\toprule
Item & Tested with & Notes \\
\midrule
Fortran compiler + MPI wrapper \code{mpif90} & gfortran 15.2 + OpenMPI
  (conda env \code{rmg-dev}) & Intel \code{ifort} + MKL also worked
  historically (settings kept as comments) \\
BLAS + LAPACK & OpenBLAS 0.3.30 & provides both; linked as
  \code{-lopenblas} \\
OpenMP & gfortran \code{-fopenmp} & \\
\bottomrule
\end{tabularx}

On the development machine, everything comes from the \code{rmg-dev} conda
environment. Activate it before you build \emph{and} before you run, because
\code{qtes.x} loads OpenBLAS and the MPI libraries from that environment at
run time.

\section{Build}

\begin{lstlisting}[style=bash]
(cd dftb && make libdftb)   # 1. DFTB library  -> dftb/libdftb.a
make                        # 2. main program  -> qtes.x
\end{lstlisting}

The top-level \code{make} does not build the library, so step~1 is required
first, and again after any change in \code{dftb/}.

\paragraph{Compiler flags.}
\begin{itemize}
\item \code{-O2 -fopenmp -fallow-argument-mismatch -std=legacy}, set in
  \code{Makefile} (\code{switch}, \code{LIB}) and in
  \code{dftb/GNUmake.inc} (\code{FOPT}).
\item \code{-fallow-argument-mismatch} is required with gfortran 10 and
  later. The legacy code passes arrays of different types and ranks to the
  same routines, which newer gfortran rejects as errors unless this flag is
  set. With it, they become warnings, and the many warnings during the
  build are expected.
\item To use Intel/MKL instead, restore the commented lines
  \code{-mkl=sequential} and \code{-openmp} in both files.
\end{itemize}

\paragraph{Clean rebuild.}
\begin{lstlisting}[style=bash]
make allclean && (cd dftb && make clean && rm -f *.o libdftb.a)
\end{lstlisting}
\code{make allclean} also deletes \code{fort.*} and \code{*.log} in the top
directory.

\section{Debug build}

To locate a floating-point error, build a separate copy with:
\begin{lstlisting}
-O0 -g -fbacktrace -ffpe-trap=zero,invalid
-fallow-argument-mismatch -std=legacy -fopenmp
\end{lstlisting}
The program then stops at the first division by zero or invalid operation
and prints a backtrace with file and line numbers.

%======================================================================
\chapter{Quick start}

\begin{lstlisting}[style=bash]
conda activate rmg-dev                     # environment with mpif90/OpenBLAS
(cd dftb && make libdftb) && make          # build
tests/run_example.sh h-c4h9-14 4           # ~1 s smoke test on 4 MPI ranks
tests/run_example.sh h-c4h9-14-ref 4       # main regression test
R=$(ls -d tests/runs/h-c4h9-14-ref_* | tail -1)
diff -q $R/ener examples/h-c4h9-14-ref/reference/ener && echo "matches reference"
\end{lstlisting}

\code{tests/run_example.sh <example> [ranks]} does the following:
\begin{itemize}
\item creates a fresh directory \code{tests/runs/<example>_<timestamp>/};
\item copies in the example's inputs, the Slater--Koster files, any restart
  checkpoint, and \code{qtes.x};
\item sets \code{OMP_NUM_THREADS=1} (unless already set) and
  \code{OPENBLAS_NUM_THREADS=1}, and raises the stack limit;
\item runs \code{mpirun -np <ranks> ./qtes.x > stdout.log}.
\end{itemize}

\textbf{Never run \code{qtes.x} inside \code{examples/}.} The program writes
its output into the current directory and would overwrite the reference
data.

%======================================================================
\chapter{Input reference}
\label{ch:input}

The program takes \textbf{no command-line arguments}. It reads these files
from the current directory, under these exact names:

\begin{tabularx}{\linewidth}{L{3.4cm} Y L{4.6cm}}
\toprule
File & Content & Read by \\
\midrule
\code{INdyn} & dynamics settings + one line per atom (wavepacket,
  displacement, momentum, mass) & \code{read_input_mpi.f} \\
\code{INxyz} & model and analysis parameters + reference geometry &
  \code{read_input_mpi.f} \\
\code{IN_dftb} & DFTB atom types and Slater--Koster file names &
  \code{dftb/dftb_init_mpi.f} \\
\code{*.spl} & Slater--Koster parameter tables & DFTB library \\
\code{traj.sav} (+ old outputs) & only for a restart &
  \code{read_input_mpi.f}, \code{set_traj.f} \\
\bottomrule
\end{tabularx}

All files are read with Fortran list-directed input. Values are separated by
spaces, and text after the last value on a line, such as the
\texttt{\textbackslash comment} used in the examples, is ignored. Lines 1--9 of \code{INdyn}
must each hold exactly the values listed below, except where values are
marked optional.

\section{\texttt{INdyn}}
\label{sec:indyn}

\needspace{8\baselineskip}
\begin{xltabular}{\linewidth}{c L{3.0cm} Y L{3.4cm}}
\toprule
Line & Values & Meaning & Units, notes \\
\midrule
\endhead
1 & \code{nat} & total number of atoms & must equal the number of geometry
  lines in \code{INxyz} \\
2 & \code{nq} & number of \textbf{quantum} atoms; they must be listed
  \textbf{first} & examples: 1 (the H) \\
3 & \code{kmax} & number of large time steps & on restart, reduced
  automatically by the time already done \\
4 & \code{kout} & substeps per large step & \\
5 & \code{dt} & large time step & a.u.\ of time; the integration step is
  \code{dt/kout} \\
6 & \code{ipot isr} & \code{ipot}: quantum potential (0 off, 1 3-D LQF,
  2 LQF along $z$). \code{isr}: initial sampling (0 off, 1 Gaussian) &
  \code{ipot > 0} requires \code{isr = 1} (ch.~\ref{ch:modes}) \\
7 & \code{pow} & sampling cutoff: keep a sampled point if its Gaussian
  weight exceeds $10^{-\mathtt{pow}}$ & used only with \code{isr = 1} \\
8 & \code{ipf} & 0 = Ehrenfest (averaged classical forces); any other value
  = independent classical copies & ch.~\ref{ch:modes} \\
9 & \code{np ist [nqtprt [nclprt]]} & \code{np}: number of QTs.
  \code{ist}: 0 = fresh start, 1 = restart from \code{traj.sav}. Optional
  print counts, default 20 each, $\le 0$ = all &
  sec.~\ref{sec:restart}, ch.~\ref{ch:output} \\
10+ & one line per atom & in the same order as \code{INxyz} (blank lines
  are skipped) & see below \\
\bottomrule
\end{xltabular}

\paragraph{Atom lines.}
\begin{lstlisting}
quantum atom:    sym  a0   x0 y0 z0   px py pz   m
classical atom:  sym  (ignored)  x0 y0 z0   (px py pz ignored)   m
\end{lstlisting}

\begin{tabularx}{\linewidth}{L{2.0cm} L{5.0cm} Y}
\toprule
Field & Meaning & Units, notes \\
\midrule
\code{sym} & element symbol (informational) & \\
\code{a0} & width parameter of the initial Gaussian, the same for $x$, $y$,
  $z$ & bohr$^{-2}$; sampling spread $\sigma = 1/\sqrt{4 a_0}$ \\
\code{x0 y0 z0} & displacement added to the \code{INxyz} position & bohr \\
\code{px py pz} & initial momentum of the wavepacket centre & a.u. Read for
  every atom, but \textbf{used only for quantum atoms}; classical atoms
  always start at rest \\
\code{m} & mass & electron masses (H = 1836, C = 22032). The
  \code{h-c4h9-14*} examples use masses $\times 10^8$ for the classical
  atoms, which effectively freezes the target \\
\bottomrule
\end{tabularx}

\needspace{12\baselineskip}
\section{\texttt{INxyz}}

\needspace{8\baselineskip}
\begin{xltabular}{\linewidth}{c L{2.2cm} L{2.2cm} Y}
\toprule
Line & Values & Used in DFTB runs? & Meaning \\
\midrule
\endhead
1 & \code{ql qls} & no & lattice constant and interaction cutoff of the old
  analytic model potentials (\code{derivs0/1/3}) \\
2 & \code{dw sc} & no & force constant and coupling (\code{sc} is
  multiplied by \code{dw}) of the model potentials \\
3 & \code{dch zh ch} & \textbf{\code{ch} only} & \code{dch}, \code{zh}:
  model-potential parameters. \code{ch}: reference $z$ (bohr) used for the
  bonding probability in \code{aver} (sec.~\ref{sec:bonding}) \\
4 & \code{esh} & \textbf{yes} & energy shift subtracted from every DFTB
  energy, $V = E_\mathrm{DFTB} - \mathtt{esh}$ (hartree) \\
5 & \code{bnd} & \textbf{yes} & threshold for the bonding probability in
  \code{aver} (bohr) \\
6 & \code{qmn qmx} & \textbf{yes} & $z$-window for the quantum atom (bohr),
  ch.~\ref{ch:modes} \\
7+ & \code{sym x y z itype} & yes & reference geometry, one line per atom,
  \textbf{in \AA ngstr\"om} (converted to bohr on input with
  0.52918~\AA/bohr). \code{itype} is the DFTB atom type, an index into
  \code{IN_dftb}. Blank lines before the geometry are allowed; lines after
  the first \code{nat} atoms are ignored \\
\bottomrule
\end{xltabular}

The program builds the actual starting positions from the \code{INxyz}
geometry plus the \code{INdyn} displacements. Make sure no two atoms end up
on top of each other: DFTB then fails at the first step with
\code{ewevge: ier = ...}. See \code{KNOWN_ISSUES.md} F2 and
ch.~\ref{ch:trouble}.

\section{\texttt{IN\_dftb}}

\begin{lstlisting}
2                ! ntype: number of atom types
2 1              ! lmax per type: number of angular shells (1 = s, 2 = s+p); type 1 = C, type 2 = H
'cc.spl'         ! SK files for every ordered pair (1,1), (1,2), (2,1), (2,2)
'ch-g2ms.spl'
'hc-g2ms.spl'
'hh.spl'
\end{lstlisting}

\begin{itemize}
\item \textbf{Slater--Koster files:} \code{examples/sk/} contains
  \code{cc.spl}, \code{hh.spl} and three C--H parameterizations:
  \code{ch/hc-g2ms.spl} (used by all examples), \code{ch/hc-b3.spl}, and
  \code{ch/hc.spl}.
\item \textbf{Hard-coded DFTB settings} in \code{derivs} (\code{qtes.f}):
  \begin{itemize}
  \item electronic temperature \code{telec} = 6000~K (Mermin free energy);
  \item neutral system (\code{qnel = 0});
  \item SCC tolerance $10^{-6}$;
  \item \code{ntype = 2}.
  \end{itemize}
  Systems with more than two element types therefore need a code change.
\end{itemize}

\section{Restart files}
\label{sec:restart}

With \code{ist = 1} on line 9:
\begin{itemize}
\item \code{traj.sav} supplies the time and all QT positions, momenta,
  actions and weights.
\item The output files \code{traj aver qt.dat pfit qtx qty qtz ener clx cly
  clz wft cor} must already exist. They are opened with
  \code{status='old'} and appended to.
\item Keep all other inputs, including the print options, the same as in
  the original run. See example \code{h-c4h9-14-restart}.
\end{itemize}

%======================================================================
\chapter{Run modes}
\label{ch:modes}

This chapter is a summary. \code{docs/DYNAMICS.md} describes each switch in
detail, with the code locations.

\needspace{8\baselineskip}
\begin{xltabular}{\linewidth}{L{2.8cm} L{2.4cm} Y L{3.6cm}}
\toprule
Mode & Switch & Typical use & Requirements \\
\midrule
\endhead
Ehrenfest & \code{ipf = 0} & classical target responds to the averaged
  quantum atom & default in all examples \\
Independent classical copies & \code{ipf} $\neq 0$ & each QT carries its
  own target & only differs from Ehrenfest when \code{isr = 1} \\
No quantum potential & \code{ipot = 0} & quantum atom moves classically
  (ensemble of classical trajectories) & --- \\
3-D LQF & \code{ipot = 1} & quantum effects (zero-point energy, spreading)
  in $x$, $y$, $z$ & \code{isr = 1} and enough QTs ($\ge 4$ with $w > 0$,
  in practice many more) \\
LQF along $z$ & \code{ipot = 2} & quantum effects along the surface normal
  only & \code{isr = 1} \\
Sampling & \code{isr = 1}, \code{pow} & wavepacket represented by randomly
  sampled QTs & --- \\
$z$-window & \code{qmn}, \code{qmx} & below \code{qmn}: QT removed from the
  ensemble; above \code{qmx}: DFTB skipped (free flight) & choose
  \code{qmx} where the potential is flat \\
Restart & \code{ist = 1} & continue a run from \code{traj.sav} &
  sec.~\ref{sec:restart} \\
Print options & \code{nqtprt}, \code{nclprt} & width of
  \code{qtx..clz} & ch.~\ref{ch:output} \\
\bottomrule
\end{xltabular}

\todo{} Physical guidance on choosing modes and parameters, for example how
many QTs, $a_0$, \code{dt}, and \code{qmn}/\code{qmx}, with the code
authors.

%======================================================================
\chapter{Running}
\label{ch:running}

\begin{lstlisting}[style=bash]
export OMP_NUM_THREADS=1 OPENBLAS_NUM_THREADS=1
ulimit -s unlimited                  # large local arrays live on the stack
mpirun -np <ranks> ./qtes.x > stdout.log 2>&1
\end{lstlisting}

\paragraph{Parallelization.}
\begin{itemize}
\item Every rank holds a full copy of all QTs.
\item The DFTB calculations are split into contiguous blocks of
  $\lceil n_p / n_\mathrm{ranks} \rceil$ QTs per rank.
\item The forces are gathered and broadcast after every substep.
\item \textbf{Rank count rule: every rank must get at least one QT}, i.e.
  \begin{equation}
    (n_\mathrm{ranks} - 1)\,\lceil n_p / n_\mathrm{ranks} \rceil < n_p .
  \end{equation}
  Otherwise the run \textbf{hangs} at the start, with no error message
  (\code{KNOWN_ISSUES.md} \#8).
\item Uneven splits are fine. For example, 48 QTs run on 4, 5, 7 or 10
  ranks with bit-identical results, but hang on 13 ranks.
\item The safest choice is a divisor of $n_p$. Results do not depend on the
  rank count.
\end{itemize}

\paragraph{Threads.} Keep \code{OPENBLAS_NUM_THREADS=1}. This matches the
original \code{-mkl=sequential} and avoids oversubscribing cores with MPI
ranks $\times$ BLAS threads.

\paragraph{Stack.} \code{ulimit -s unlimited} is needed. With
\code{-fopenmp}, large local arrays (about 28~MB in \code{trot}) are
allocated on the stack.

\paragraph{Cluster.} An old SLURM example is in \code{docs/sbatch.script}
(24 tasks per node, \code{OMP_NUM_THREADS=1}).

\paragraph{Measured run times} (gfortran + OpenBLAS, development machine):

\begin{tabularx}{\linewidth}{Y r r r r r}
\toprule
Example & Atoms & QTs & Substeps & Ranks & Wall time \\
\midrule
\code{h-c4h9-14} & 14 & 48 & 8 & 4 & $\sim$1 s \\
\code{h-c4h9-14-ref} & 14 & 48 & 50 & 4 & $\sim$3 s \\
\code{h-flake-53} & 53 & 48 & 30 & 48 & $\sim$4 s \\
\code{h-flake-111} & 111 & 12 & 50 & 12 & $\sim$25 s \\
\bottomrule
\end{tabularx}

%======================================================================
\chapter{Output reference}
\label{ch:output}

All output files are written by MPI rank 0 into the current directory.
\begin{itemize}
\item \textbf{Units:} time in a.u., positions in bohr, energies in hartree.
\item \textbf{Columns:} in all files except \code{wf0}, \code{wft},
  \code{qt.dat} and \code{traj.sav}, column 1 is the time.
\item \textbf{When written:}
  \begin{itemize}
  \item per substep: \code{ener}, \code{cor};
  \item per large step: \code{aver}, \code{traj}, \code{qtx..clz},
    \code{pfit}, plus $t = 0$ for most of them;
  \item \code{qt.dat}: per large step, without $t = 0$;
  \item \code{traj.sav}: overwritten at every large step.
  \end{itemize}
\end{itemize}

\section{Files}

\needspace{8\baselineskip}
\begin{xltabular}{\linewidth}{L{3.0cm} Y}
\toprule
File & Content and columns \\
\midrule
\endhead
\code{ener} & $t$, $E_\mathrm{total}$, $\langle Q\rangle$, $\langle
  V\rangle$, $\langle K_\mathrm{quantum}\rangle$, $\langle
  K_\mathrm{classical}\rangle$: weighted averages over the QTs; $V$
  includes the \code{esh} shift. Energy conservation of $E_\mathrm{total}$
  is the main health check \\
\code{cor} & $2t$, $\mathrm{Re}\,C$, $\mathrm{Im}\,C$: correlation function
  $C = \sum_i w_i e^{2 i s_i} / \sum_i w_i$ (an autocorrelation if
  $\psi(0)$ is real) \\
\code{aver} & $t$, norm, $P$, $\langle x_1\rangle \ldots \langle
  x_{40}\rangle$: norm $= \sum_i w_i$ (drops below 1 when QTs leave below
  \code{qmn}); $P$ = bonding probability (sec.~\ref{sec:bonding});
  weighted average of the first 40 coordinates \\
\code{traj} & $t$, the first $4\,n_{qd}$ coordinates of QT 2, and $s_2$: a
  representative trajectory \\
\code{qtx qty qtz} & $t$, then $x$/$y$/$z$ of the quantum atom in QT 1
  \ldots\ $\min(n_p, \mathtt{nqtprt})$ \\
\code{clx cly clz} & \code{ipf = 0}: $t$, then $x$/$y$/$z$ of classical
  atoms 1 \ldots\ $\min(n_\mathrm{cl}, \mathtt{nclprt})$ from the first QT
  with $w > 0$. \code{ipf} $\neq 0$: the first classical atom in QT 1
  \ldots\ $\min(n_p, \mathtt{nqtprt})$ \\
\code{pfit} & $t$ and the LQF fit coefficients (\code{ipot > 0}; empty
  otherwise) \\
\code{qt.dat} & 2 lines per QT per large step: all positions + $w$, then
  all momenta + $s$ (full precision) \\
\code{traj.sav} & restart checkpoint: time, then the same 2 lines per QT \\
\code{wf0}, \code{wft} & initial and final positions of every QT (first 40
  coordinates). On restart, \code{wf0} holds the restart state and
  \code{wft} is appended to \\
\code{trajt}, \code{dftbt} & timings: start, end, elapsed. They differ on
  every run \\
\code{out.dftb} & echo of the DFTB setup read from \code{IN_dftb} \\
\code{stdout.log} & screen output: setup summary, including
  \code{PRINT OPTIONS}, initial energies, step counter \\
\code{fort.1000+rank} & per-rank trace written at initialization \\
\bottomrule
\end{xltabular}

\code{docs/sg-notes-files} gives the column formats in the original
notation.

\section{Bonding probability \texorpdfstring{$P$}{P} (column 3 of
\texttt{aver})}
\label{sec:bonding}

\begin{equation}
  P = \frac{\sum_i w_i\, \Theta\!\left(|z_{\mathrm{H},i} - \mathtt{ch}|
      - \mathtt{bnd}\right)}{\sum_i w_i},
\end{equation}
where $\Theta$ is the step function ($\Theta(x) = 1$ for $x \ge 0$),
$z_{\mathrm{H},i}$ is the \textbf{absolute $z$ coordinate} of the quantum
atom in QT $i$, and \code{ch}, \code{bnd} come from \code{INxyz}.

With the values in the current examples (\code{ch} = 0.65 bohr, \code{bnd}
= 1 bohr), the H is always more than \code{bnd} away from $z = \mathtt{ch}$,
so $P = 1$ throughout. \todo[to be clarified with the authors] What $P$ is
meant to measure. A commented-out older line used the H--C distance instead
of the absolute $z$.

%======================================================================
\chapter{Examples}

All examples are in \code{examples/}. \code{examples/README.md} describes
each one in detail, including how its reference data was made.

\begin{tabularx}{\linewidth}{L{3.6cm} Y L{3.6cm}}
\toprule
Example & System & Purpose \\
\midrule
\code{h-c4h9-14-ref} & H + C$_4$H$_9$ (frozen), 48 QTs & main regression
  test \\
\code{h-c4h9-14-restart} & same, restart from $t = 80$ & tests the restart
  path \\
\code{h-c4h9-14} & same, short & $\sim$1 s smoke test \\
\code{h-flake-53} & H + 52-atom graphene flake & larger system, Ehrenfest
  target \\
\code{h-flake-53-print} & same, \code{nqtprt = 48}, \code{nclprt = 0} &
  demonstrates the print options \\
\code{h-flake-111} & H chemisorbed on a 110-atom flake, \code{ipot = 1},
  \code{isr = 1} & quantum potential with sampling \\
\bottomrule
\end{tabularx}

The graphene-flake examples correspond to the system studied in
ref.~\cite{wang2014}.

\todo{} Step-by-step tutorials with plots.

%======================================================================
\chapter{Verification and regression tests}

\begin{itemize}
\item \textbf{Reproducibility:} runs are bit-reproducible, and the results
  are the same for any valid rank count. This was tested with 4, 5, 7 and
  10 ranks on \code{h-c4h9-14-ref}. Every example has a \code{reference/}
  folder.
\item \textbf{After a code change:} run all examples. Every file except
  \code{trajt}/\code{dftbt} must be byte-identical to the reference, unless
  the change is meant to alter results.
\end{itemize}

\begin{lstlisting}[style=bash]
for ex in "h-c4h9-14-ref 4" "h-c4h9-14-restart 4" "h-flake-53 48" \
          "h-flake-53-print 48" "h-flake-111 12"; do
  set -- $ex; tests/run_example.sh $1 $2 >/dev/null
  R=$(ls -d tests/runs/$1_2* | tail -1)
  for f in ener cor aver qt.dat traj.sav qtx clx; do
    cmp -s $R/$f examples/$1/reference/$f || echo "$1: $f differs"
  done
done
\end{lstlisting}

\paragraph{Physics checks.}
\begin{itemize}
\item $E_\mathrm{total}$ in \code{ener} should be conserved, to all printed
  digits in the examples.
\item A restart from $t = 80$ reproduces a single run from 0 to 100 to
  round-off.
\item With per-rank checksums, all ranks stay bit-identical (see
  \code{KNOWN_ISSUES.md} F7).
\end{itemize}

%======================================================================
\chapter{Troubleshooting}
\label{ch:trouble}

\needspace{8\baselineskip}
\begin{xltabular}{\linewidth}{L{4.4cm} Y L{3.8cm}}
\toprule
Message or symptom & Cause & Fix \\
\midrule
\endhead
\code{Bad integer for item 2 in list input} at \code{read_input_mpi.f} &
  a line of \code{INdyn} has too few values, e.g.\ line 6 needs
  \code{ipot isr} & give all required values (sec.~\ref{sec:indyn}) \\
\code{ewevge: ier = ...} in \code{fort.10x}, run stops before $t = 0$ &
  DFTB failed, usually because two atoms overlap after the \code{INdyn}
  displacements & check the starting distances (\code{INxyz} +
  displacements) \\
\code{QPOT ERROR ... only N QTs with w>0; fit needs M} & too few QTs left
  inside the \code{qmn} window for the quantum-force fit & more QTs, a
  lower \code{qmn}, or \code{ipot = 0} \\
\code{QPOT ERROR ... DPOSV FAILED ... (isr=0?)} & the fit matrix is
  singular: all QTs identical (\code{isr = 0}) or lying in a plane & use
  \code{isr = 1} with \code{ipot > 0} \\
energies grow quickly before a \code{QPOT ERROR} & the fit is
  ill-conditioned with few, clustered QTs & more QTs; keep them inside the
  window \\
restart stops opening a file & restart needs the previous output files &
  copy them into the run directory (sec.~\ref{sec:restart}) \\
\code{IEEE_DENORMAL} note at exit & harmless: products of near-zero
  coordinates & none \\
\code{MPI_FINALIZE ERR 0} (one per rank) & the return code of
  \code{MPI_Finalize}; 0 means success & none \\
step in $E_\mathrm{total}$ when a QT re-enters below \code{qmx} & free
  flight with a frozen potential above \code{qmx} & raise \code{qmx} to
  where the potential is flat \\
run hangs right after \code{FINISHED SAMPLING}, all ranks at 100\,\% CPU &
  some MPI rank has no QTs (\code{KNOWN_ISSUES.md} \#8) & use a rank count
  with every rank getting $\ge 1$ QT, e.g.\ a divisor of $n_p$ \\
run hangs or crashes at start with a large system & stack too small &
  \code{ulimit -s unlimited} \\
\bottomrule
\end{xltabular}

Open and fixed bugs are tracked in \code{docs/KNOWN_ISSUES.md}.

%======================================================================
\chapter{Developer notes}

\todo{} Planned content:
\begin{itemize}
\item the source layout (\code{qtes.f} subroutines \code{init_traj},
  \code{trot}, \code{qpot}, \code{over}, \code{aver}, \code{derivs};
  \code{read_input_mpi.f}; \code{set_traj.f}; \code{dftb/});
\item the replicated-data MPI design and its rules (all ranks must make
  identical decisions; see F7);
\item the fixed-form Fortran 72-column limit;
\item where to add a new potential (\code{derivs0..3}) or an output;
\item the test procedure before committing.
\end{itemize}

%======================================================================
\appendix

\chapter{Units}

\begin{tabularx}{\linewidth}{L{2.6cm} L{3.6cm} Y}
\toprule
Quantity & Unit & Conversion \\
\midrule
length & bohr ($a_0$) & 1 bohr = 0.52918~\AA\ (value used in the code) \\
time & a.u. & 1 a.u.\ = 0.024189 fs \\
energy & hartree ($E_h$) & 1 $E_h$ = 27.211 eV \\
mass & electron mass ($m_e$) & H = 1836, C = 22032 ($12 \times 1836$) \\
momentum & a.u. & $\hbar/a_0$ \\
\bottomrule
\end{tabularx}

\chapter{Glossary}

\begin{tabularx}{\linewidth}{L{3.0cm} Y}
\toprule
Term & Meaning \\
\midrule
QT & quantum trajectory: one member of the ensemble that samples the
  quantum nuclei \\
QDOF / \code{nqd} & quantum degrees of freedom: 3 $\times$ number of
  quantum atoms \\
LQF & linearized quantum force: the quantum potential from a linear fit
  over the QTs \\
Ehrenfest & the classical atoms move in the QT-averaged force
  (\code{ipf = 0}) \\
$w$, $s$ & weight and action (phase) of a QT \\
$z$-window & \code{qmn}/\code{qmx}: the $z$-range of the quantum atom inside
  which DFTB is computed \\
\bottomrule
\end{tabularx}

\chapter{Compile-time limits}

\begin{tabularx}{\linewidth}{L{2.0cm} L{3.4cm} r Y}
\toprule
Parameter & File & Value & Limits \\
\midrule
\code{ndx} & \code{bohm.inc} & 700 & total coordinates, i.e.\ at most 233
  atoms \\
\code{npx} & \code{bohm.inc} & 5000 & QTs \\
\code{nqx} & \code{bohm.inc} & 200 & quantum coordinates \\
\code{nndim} & \code{dftb/maxima.inc} & 800 & atoms in DFTB \\
\code{mdim} & \code{dftb/maxima.inc} & 3200 & DFTB orbitals \\
\code{maxtyp} & \code{dftb/maxima.inc} & 6 & DFTB atom types
  (\code{derivs} currently passes 2) \\
\bottomrule
\end{tabularx}

After changing a limit, do a clean rebuild of both \code{dftb/} and the
main program.

%======================================================================
\begin{thebibliography}{9}

\bibitem{garashchuk2013}
S.~Garashchuk, J.~Jakowski, L.~Wang, B.~G.~Sumpter,
``Quantum Trajectory--Electronic Structure Approach for Exploring Nuclear
Effects in the Dynamics of Nanomaterials'',
\emph{J. Chem. Theory Comput.} \textbf{9}, 5221--5235 (2013).
\href{https://doi.org/10.1021/ct4006147}{doi:10.1021/ct4006147}

\bibitem{wang2014}
L.~Wang, J.~Jakowski, S.~Garashchuk,
``Adsorption of a Hydrogen Atom on a Graphene Flake Examined with Quantum
Trajectory/Electronic Structure Dynamics'',
\emph{J. Phys. Chem. C} \textbf{118}, 16175--16187 (2014).
\href{https://doi.org/10.1021/jp503261k}{doi:10.1021/jp503261k}

\bibitem{wang2016}
L.~Wang, J.~Jakowski, S.~Garashchuk, B.~G.~Sumpter,
\emph{J. Chem. Theory Comput.} \textbf{12}, 4487--4500 (2016).
\href{https://doi.org/10.1021/acs.jctc.6b00126}{doi:10.1021/acs.jctc.6b00126}

\end{thebibliography}

\end{document}
